\chapter{关键原语实现与工程优化}

公共协议只看到HIBE与IBS接口，实际计算则由各后端完成。本章先说明两种HIBE和主要IBS的实现，
再展开结构化采样、紧凑编码和接收会话三项优化。每一节都从该原语在IB-ME中的职责出发，随后
给出关键方程和软件中的对象表示。

\section{两种分层身份基加密}

\subsection{LATTE：在NTRU结构中完成两层委派}

LATTE在负循环环$R_q=\mathbb Z_q[x]/(x^N+1)$上建立NTRU结构陷门
~\cite{zhao2024latte}。作品采用LATTE-3参数：$N=1024$、$q=68718428161$，分层深度为2。
建立阶段生成短多项式$f,g$，求解NTRU方程
\[
 fG-gF=q,
 \qquad h=gf^{-1}\pmod q,
 \qquad S_0=\begin{pmatrix}g&-f\\G&-F\end{pmatrix}.
\]
$h$进入公开参数，$S_0$作为根陷门保存在主状态中。

接收身份$r$经过独立域标识的哈希得到$A_1=H_1(r)$。发钥算法利用根陷门为第一层身份生成
短基，得到接收密钥$\mathsf{dk}_r$。接收者指定发送身份$s$后，再计算
$A'_2=H_2(r,s)$并委派叶钥$(t_0,t_1,t_2)$。叶钥满足
\[
 t_0+t_1h+t_2A_1=A'_2\pmod q.
\]
这条关系既用于解密，也是密钥从磁盘恢复时的绑定检查。若把另一个接收身份或发送身份放入
同一对象，等式右端随之变化，加载器会拒绝该组合。

加密算法取短噪声，把32字节种子编码进环元素，并沿$(r,s)$路径生成胶囊。解密使用叶钥消去
公开环项，恢复种子后再执行一次确定性重加密核对。该核对用于排除落在解码区间内的错误结果，
随后才把种子交给混合加密层。

实现复用了LATTE作者代码中的Gaussian采样和ffSampling路径，并把C对象包装为严格的Python
后端。接入过程中修正了COSAC读取Box--Muller新缓冲的位置；每次原生采样完成后，Python层
使用精确环算术检查NTRU关系、委派关系和对象形状。公开参数、主陷门、一级密钥、叶密钥和
密文分别使用独立对象类型，参数快照随对象保存。

LATTE的建立和一级发钥时间在三组代表中都较低，因此A、B共用它作为HIBE。相同HIBE也使两组
差异主要来自IBS，便于分析认证转换与原生模格签名在完整协议中的开销。

\subsection{Lai：逐层追加身份块}

Lai方案使用理想格陷门组织分层委派~\cite{lai2015trapdoors}。根公开向量与短陷门满足
\[
 \mathbf a_0^{\mathsf T}R+\mathbf a_1=\mathbf g.
\]
两层身份分别映射为二进制环元素$\mathrm{id}_1$和$\mathrm{id}_2$。第一层公开项与完整身份公开项
写成
\[
 \mathbf a_{\mathrm{id}_1}=\mathbf u_0+\mathrm{id}_1\mathbf u_1,
 \qquad
 \mathbf a_{\mathrm{ID}}=\mathbf a_{\mathrm{id}_1}+\mathrm{id}_2\mathbf u_2.
\]
发钥和委派每进行一层，就向公开矩阵追加对应身份块，并为扩展后的矩阵生成新陷门。接收者$r$
对应第一层，期望发送者$s$对应第二层。

加密时取短矩阵$R_0,R_1,R_2$，组合为
\[
 R_{\mathrm{ID}}=R_0+\mathrm{id}_1R_1+\mathrm{id}_2R_2,
 \qquad \mathbf e_1=-\mathbf e_0^{\mathsf T}R_{\mathrm{ID}},
\]
再计算
\[
 \mathbf u=F_1\mathbf s+p\mathbf e_0,
 \qquad
 \mathbf v=(\mathbf a_{\mathrm{ID}}+h'\mathbf g)\mathbf s
          +p\mathbf e_1+\mathbf m.
\]
叶陷门对$(\mathbf u,\mathbf v)$执行Babai反演，检查误差的模$p$结构并恢复32字节种子。

C组采用$n=2048$、$q=274877976577$、$k=39$、深度2、$p=2$的C2配置，噪声参数为
$\sigma=2$、$\sigma_r=1$。身份映射使用二进制系数，避免身份乘短矩阵时放大噪声。实现按层
追加身份块，统一矩阵形状，并用Klein/GPV路径完成陪集采样。公开参数、根陷门、两级密钥与
密文均按5字节系数编码。

Lai的主要成本在二级委派。一次完整C组解密中，委派占总时间的绝大部分；第~\ref{sec:lai-session}
节因此把固定身份对的委派移动到接收会话准备阶段。

\section{身份基签名后端}

\subsection{BNN认证转换：接入ML-DSA与Falcon}

BNN转换把普通数字签名变为身份基签名~\cite{galindo2006ibs}。系统建立时，发钥机构生成一套
机构签名密钥。为身份$s$发钥时，先为用户生成普通签名密钥对
$(\mathsf{vk}_s,\mathsf{xk}_s)$，再由机构签署
\[
 \mathsf{cert}_s=\mathsf{DS.Sign}_{\mathsf{sk}_{\rm CA}}
       (\mathsf{Enc}(s,\mathsf{vk}_s)).
\]
用户签名密钥由$\mathsf{xk}_s$、$\mathsf{vk}_s$和$\mathsf{cert}_s$组成。签消息时，用户对
上下文、身份、公钥、证书和协议消息的规范编码签名。完整IBS签名携带用户公钥、证书和用户
签名；验证者先检查证书，再检查用户签名。

同一转换核心接入ML-DSA与Falcon两个普通签名引擎。ML-DSA实现固定FIPS~204的参数与对象格式
~\cite{nist2024fips204}；A组选用ML-DSA-44，其基础公钥为1,312字节、基础签名为2,420字节，
组合后的IBS签名为6,237字节。Falcon后端采用Falcon-512和Falcon-1024的固定作者实现与
padded格式~\cite{falcon2020}。两者共享机构证书流程，但底层密钥、签名编码和版本号完全分开。

认证转换的优势是建立和在线签名较快，也便于复用标准化实现。其额外空间来自用户公钥、机构
证书和用户签名三部分。A组的实验说明，即使采用较小的ML-DSA参数，完整密文仍会受到一次性
认证和HIBE胶囊的固定开销影响。

\subsection{Luo：模格上的原生身份签名}

Luo后端实现论文的模块格IBS~\cite{luo2024ibs}。身份哈希、消息挑战、响应和提示信息直接进入
签名方程，验证者只需公开参数、身份和消息即可验签。实现保留参考格式，同时提供无损紧凑格式；
两种格式加载后得到同一数学对象，签名算法和验签方程不变。

Luo后端在组合矩阵中分别与LATTE、Lai配对。它不是A、B、C三组的重点代表，但承担两个作用：
一是验证公共协议能够接入另一种原生模格IBS；二是为十组同机补测提供与SRGG不同的模格实现。
后端公开身份上限、签名输入上限和签名尺寸，装配层据此计算应用容量。

\subsection{SRGG：身份目标与消息挑战进入同一关系}

SRGG采用模块格结构，公开矩阵写成
\[
 A=[I_d\mid\widehat A\mid B]
\]
~\cite{jia2025ibs}。身份$s$决定目标$h=H_1(P,s)$，用户短见证$y$满足
$Ay=h\pmod Q$。签名构造矩阵和响应为
\[
 M_h=[\,2A\mid Qj-2h\,]\pmod{2Q},
 \qquad
 z=p_0+\begin{bmatrix}y\\1\end{bmatrix}(\zeta u+c),
\]
其中$\zeta=1+x^{n/2}$、$j=\zeta^\star e_1$。扰动$p_0$的承诺和消息共同决定挑战$c$。
验证者由
\[
 W'=M_hz-Qjc
\]
恢复承诺，重新计算挑战，并检查响应范数。

B组采用PARAM-I：多项式次数$n=512$、模块秩2、$p=16383$、$Q=15744063$，gadget基为31、
位数为2。建立阶段生成公开矩阵和短陷门；身份发钥执行相关扰动与短原像采样；签名阶段重新
采样承诺和挑战。实现同时保留JSON参考格式与紧凑格式，二者由同一后端运算。第
~\ref{sec:srgg-codec}节说明四类对象的位级编码。

\subsection{Sageloli：环陷门与非齐次签名方程}

Sageloli的公开向量与主陷门$R$满足~\cite{sageloli2023ibs}
\[
 A=[\bar A\mid g-\bar AR].
\]
为身份$s$发钥时，算法先取身份盐$\rho$，计算$h=H_1(P,\rho,s)$，再采样身份陷门$T_s$，使
\[
 AT_s=h-g\pmod q.
\]
签名另取消息盐$\eta$，通过相关扰动和三进制陪集采样得到响应$z$。验证者检查范数及
\[
 [A,h]z=H_2(P,\rho,s,\eta,M')\pmod q.
\]

C组采用PARAM-I：$(d,k,\ell)=(2048,65,132)$、$q=3^{65}$、$r^2=49$；身份与签名采样宽度
平方分别为$2720281532542$和$270116709797628319132317$。PARAM-II保持同一算法而增大参数。
两档均使用环原生后端、规范对象和固定身份发钥服务。

该方案的难点不是验签方程，而是论文参数下的相关扰动。若把环乘法展开为普通矩阵，单个中间
矩阵就超过常规工作站内存。下一节说明作品如何保留环结构完成采样。

\subsection{Wang--PW：目标规模矩阵SIS实现}

这里的Wang--PW指Wang基本IBS加Pan--Wagner自适应转换，它与第~2章所述的Wang通用IB-ME
构造是两个层次：前者是可替换的IBS后端，后者规定HIBE与IBS怎样组合。基本IBS采用矩阵SIS
路线，Pan--Wagner转换为身份和消息分别加入随机盐~\cite{wang2024ibme,pan2021ibs}。

目标档参数为
\[
 n=512,\qquad q=4294967291,\qquad s_0=4.
\]
令$k=\lceil\log_2q\rceil=32$，则gadget宽度$w=nk=16384$。实现取
$\bar m=w$、$m=\bar m+w=32768$，签名向量维数为
$M=m+w=49152$。根公开矩阵由G陷门构造：
\[
 A=[\bar A\mid G-\bar AR],
\]
其中短矩阵$R$只保存在发行机构的主状态中。

发钥时，Pan--Wagner外层先为身份$s$生成固定身份盐$\rho$，再计算
$H_1(P,\rho,s)$。后端逐列采样用户陷门$X$，使
\[
 AX+H_1=G,
 \qquad
 F_s=[A\mid H_1],
 \qquad
 F_s\begin{bmatrix}X\\I\end{bmatrix}=G\pmod q.
\]
用户密钥只保存$X$、身份盐和公开上下文，不包含主陷门$R$。签名时重新生成消息盐$\eta$，
计算消息目标$H_2(P,\rho,s,\eta,M')$，再使用用户G陷门采样短向量$z$：
\[
 F_s z=H_2(P,\rho,s,\eta,M')\pmod q.
\]
完整签名为$(\rho,\eta,z)$。验签重新计算$H_1$和$H_2$，检查维数、非零条件、整数平方范数和
模方程。

\begin{figure}[H]
\centering
\begin{tikzpicture}[node distance=7mm and 7mm]
 \node[flow,text width=35mm] (root) {Setup\\生成$A$与主陷门$R$};
 \node[flow,right=of root,text width=36mm] (key) {KeyExtract\\身份盐$\rho$与用户陷门$X$};
 \node[flow,right=of key,text width=36mm] (sign) {Sign\\消息盐$\eta$与短原像$z$};
 \node[check,below=of key,text width=43mm] (verify) {Verify\\公开重算目标、范数和模方程};
 \draw[flowarrow](root)--(key);
 \draw[flowarrow](key)--(sign);
 \draw[flowarrow](sign)--(verify);
\end{tikzpicture}
\caption{Wang--PW后端的身份发钥、签名与公开验签}
\label{fig:wang-pw-flow}
\end{figure}

有限离散Gaussian由独立C++组件执行。实现使用256 bit MPFR精度、192 bit概率质量和20个标准差
截断窗口；根采样、用户发钥和签名分别使用固定宽度与公开门限。Python层负责参数快照、域分隔
哈希、对象编码和精确模方程，原生组件只承担有限Gaussian采样。

目标档的公钥约64.0~MiB，主状态约576~MiB，用户签名密钥约2.06~GiB，签名约384~KiB。
大对象采用流式保存和恢复，签名进程只读取用户密钥
与公开参数。固定发行服务记录身份盐和已发布密钥，使同一身份在重启后继续使用同一用户对象。

该实现已完成独立签名、公开验签、错误身份与消息、签名篡改、状态恢复以及LATTE组合。项目按
约115 bit量子安全工作值使用该目标档。它的Setup、发钥和签名消耗显著高于常规十组，第~6章
单独列出分阶段实测结果。

\section{利用环结构降低采样开销}
\subsection{大矩阵从哪里产生}
Sageloli的发钥和签名都需要在格陪集中采样。相关扰动由陷门决定其形状，再与离散采样配合，得到满足方程和范数要求的短向量。实现既要得到正确的相关性，也要承受论文参数下的对象规模。

环元素是一组多项式系数，乘法可以直接在负循环环上计算。如果把每个环元素展开成旋转矩阵，后续陷门与协方差也会随之变成很大的实数矩阵。PARAM-I取$d=2048$、$\ell=132$、$k=65$，签名相关扰动的系数维数为
\[
 D=d(\ell+k)=403456.
\]
仅显式保存一个$D\times D$实数矩阵，按每项8字节估算就需要约1.18 TiB，还未计入分解工作区。因此，实际实现需要利用原式的结构来组织计算。

本作品始终保存多项式系数。公钥、主陷门和身份陷门分别以$\ell\times d$、$(\ell-k)\times k\times d$和$\ell\times k\times d$的系数数组表示。环乘法在原环上执行，结果仍检查
\[
 A=[\bar A\mid g-\bar AR],\qquad AT_s=H_1(P,\rho,s)-g\pmod q.
\]
这样，公钥和密钥关系保持不变，计算过程避免展开旋转矩阵。

\subsection{把大分解转为逐频率的小分解}
设当前陷门为$T$，其形状为$\ell\times k$个环元素，令
\[
 J=\begin{bmatrix}-T\\I_k\end{bmatrix},\qquad
 a=s^2-r^2,\qquad c=9r^2.
\]
这里$r$是基础采样宽度，$s$是该次相关扰动使用的宽度。需要计算的协方差形状为
\[
 \Sigma=aI-cJJ^\star,
\]
其中$\star$表示环上的伴随。这个式子的主要部分是标量矩阵$aI$，修正部分由只有$k$列的$J$产生。利用这一点，可以将分解限制在较小的子空间。

第一步，对多项式执行负循环Fourier变换。在每个频率$\nu$上，环乘法变为复数乘法，$J$对应一个$(\ell+k)\times k$的复数矩阵$J_\nu$。第二步，计算Gram矩阵
\[
 G_\nu=J_\nu^*J_\nu
      =U_\nu\operatorname{diag}(\lambda_i)U_\nu^*.
\]
$G_\nu$只有$k\times k$；在PARAM-I中，最大特征分解为$65\times65$。

\begin{figure}[H]
\centering
\begin{tikzpicture}[node distance=8mm and 7mm]
 \node[flow,text width=39mm] (poly) {陷门多项式$T$\\保存系数数组};
 \node[flow,text width=39mm,right=of poly] (fourier) {负循环Fourier变换\\$J_\nu$：$197\times65$};
 \node[flow,text width=39mm,right=of fourier] (gram) {计算$J_\nu^*J_\nu$\\分解$65\times65$矩阵};
 \node[check,text width=39mm,below=of gram] (apply) {对新随机向量应用\\低秩平方根};
 \node[check,text width=39mm,left=of apply] (inverse) {逆变换回系数域\\逐系数离散采样};
 \node[check,text width=39mm,left=of inverse] (check) {检查精确模方程\\和原方案范数条件};
 \draw[flowarrow](poly)--(fourier);
 \draw[flowarrow](fourier)--(gram);
 \draw[flowarrow](gram)--(apply);
 \draw[flowarrow](apply)--(inverse);
 \draw[flowarrow](inverse)--(check);
\end{tikzpicture}
\caption{结构化采样的计算步骤；矩阵形状对应PARAM-I}
\label{fig:ring-sampler}
\end{figure}

第三步，利用小矩阵的特征分解构造平方根算子
\[
 L_\nu=\sqrt a I-
 J_\nu U_\nu\operatorname{diag}\!\left(
 \frac{c}{\sqrt a+\sqrt{a-c\lambda_i}}
 \right)U_\nu^*J_\nu^*.
\]
该式满足$L_\nu L_\nu^*=aI-cJ_\nu J_\nu^*$。对新随机向量应用$L_\nu$时，按右侧的小矩阵与向量乘法逐步计算，无需形成完整的$L_\nu$。随后逆变换回系数域，继续执行原方案要求的离散Gaussian和三进制陪集采样。方法使用格陷门采样中的扰动思想~\cite{micciancio2012trapdoors}，并针对本方案的环结构组织执行。

需要保存的主要数据量为$O(d\ell k+dk^2)$，对应陷门系数和各频率的小矩阵计划。相较$O(D^2)$的稠密表示，这直接降低了表示和分解规模。已经计算好的计划可在同一陷门下复用，每次发钥和签名仍生成新的随机样本。

\subsection{数值计算与精确方程怎样衔接}
Fourier步骤使用数值运算，方案中的模方程则必须精确成立。实现将两部分分开：数值计划在建立时检查正定充分条件和变换误差；回到系数域后，环方程由整数运算检查。大模数$q=3^k$的公开乘积使用FLINT整数多项式，私有小整数乘法采用双素数NTT与CRT恢复，并检查整数卷积幅度落在可唯一恢复的范围内。

签名生成后继续执行原方案的范数门与验签方程。这样，浮点计算的职责是产生相关扰动，精确算术的职责是验证代数关系；两者在采样接口处衔接。第6章给出PARAM-I与PARAM-II的实际运行成本。

\section{按取值范围编码密钥与签名}
\label{sec:srgg-codec}
\subsection{JSON表示中的重复开销}
SRGG的公钥、陷门和签名含有大量多项式系数。JSON将每个整数写成十进制字符串，还需要负号、逗号、括号和字段名。对大对象，表示开销会累积到与数学数据相近甚至更大的规模。

紧凑编码依据字段的数学取值范围选择位数。以取值为$-1,0,1$的陷门系数为例，三种可能值只需2 bit；公开模$Q=15744063$的系数落在$[0,Q)$，使用24 bit。对短密钥和签名响应，先以Zigzag映射处理正负号，再选取能够容纳该字段的最小统一位宽。

\begin{table}[H]
\centering\small
\caption{SRGG不同字段的编码方法}
\label{tab:codec-fields}
\begin{tabularx}{\textwidth}{@{}p{.23\textwidth}p{.25\textwidth}X@{}}
\toprule
\rowcolor{black!8} 数据 & 存储方法 & 解码时检查 \\
\midrule
公开模$Q$系数 & 每系数24 bit & $0\leq x<Q$ \\
三值陷门系数 & 每系数2 bit & 拒绝第四个未使用码字 \\
短密钥、响应整数 & Zigzag与最小统一位宽 & 数值范围、位宽最小性 \\
二元挑战 & 每系数1 bit & 固定多项式次数与长度 \\
类型、参数与上下文 & 固定头部和长度前缀 & 对象类别、公钥、参数一致 \\
\bottomrule
\end{tabularx}
\end{table}

对含$N$个系数、统一位宽为$b$的字段，其主要编码长度为$\lceil Nb/8\rceil$字节，再加必要头部。解码后得到与JSON格式相同的数学对象，签名和验证继续使用原算法。

\subsection{四类对象与规范解码}
该方法同时用于公钥、主状态、用户密钥和签名。各类对象有独立类型标识，系数顺序固定，末字节未使用位必须为零。解析器先根据形状计算长度，再读取数据，并检查参数、身份、公钥上下文及取值范围。解码后重新编码应得到相同字节串，从而保证同一个对象只有一种规范表示。

对象尺寸的改善会减少密钥文件和传输字节，实际收益取决于该对象在完整协议中的占比。例如，签名变短可以直接缩短对称加密载荷，但HIBE胶囊和一次性认证仍占固定空间。第6章将四类对象的编码对照与完整密文组成放在一起分析。编解码也有计算成本，实验同时保留该项测量。

\section{通过接收会话减少重复委派}
\label{sec:lai-session}
\subsection{重复计算发生在哪里}
在Lai组合中，接收者持有第一层密钥$\mathsf{dk}_r$。解密来自发送者$s$的消息前，需要运行
\[
 \mathsf{dk}_{r,s}\leftarrow\mathsf{Derive}(\mathsf{dk}_r,s).
\]
同一接收者连续读取同一发送者的消息时，委派输入保持相同；反复执行会重复支付生成叶钥的开销。加密种子、密文与每条消息的签名则各不相同，需要逐条处理。

本作品引入显式接收会话：在会话建立时生成并验证一次叶钥，后续在固定双方身份下复用。图~\ref{fig:receiver-session}展示计算安排的变化。

\begin{figure}[H]
\centering
\begin{tikzpicture}[node distance=6mm and 6mm]
 \node[font=\small,anchor=east] (old) {逐条委派};
 \node[flow,right=of old,text width=30mm] (o1) {委派＋解密1};
 \node[flow,right=of o1,text width=30mm] (o2) {委派＋解密2};
 \node[flow,right=of o2,text width=30mm] (o3) {委派＋解密3};
 \node[font=\small,anchor=east,below=22mm of old] (new) {接收会话};
 \node[check,right=of new,text width=30mm] (prep) {委派一次\\验证并保留叶钥};
 \node[check,right=of prep,text width=30mm] (dec) {解密1、2、3\\逐条认证与验签};
 \node[check,right=of dec,text width=30mm] (close) {关闭会话\\停止使用叶钥};
 \draw[flowarrow](o1)--(o2);\draw[flowarrow](o2)--(o3);
 \draw[flowarrow](prep)--(dec);\draw[flowarrow](dec)--(close);
\end{tikzpicture}
\caption{会话把与密文无关的委派安排到准备阶段}
\label{fig:receiver-session}
\end{figure}

\subsection{身份绑定与会话生命周期}
会话建立时检查公开上下文、接收者身份和一级钥绑定；生成叶钥后检查完整$(r,s)$路径及陷门方程。会话只服务于该公钥和身份对，关闭后停止解密。每条消息继续执行规范解析、外层字节认证、HIBE解封装、GCM验证、IBS验签及一次性公钥承诺核对。

叶钥作为秘密状态保留在进程内，使用者显式关闭会话，会话对象不支持序列化。长期保存的接收密钥仍采用原持久格式。该设计适合双方身份固定的连续通信，并带来明确的额外密钥存储成本。

\subsection{怎样计算连续消息的平均成本}
设一次会话准备耗时为$T_{\rm prep}$，会话内单条解密为$T_{\rm sess}$。处理$N$条消息的总成本与平均成本分别为
\[
 T_{\rm total}(N)=T_{\rm prep}+NT_{\rm sess},\qquad
 \bar T(N)=\frac{T_{\rm prep}}{N}+T_{\rm sess}.
\]
首条消息承担全部准备成本；连续消息逐渐摊薄这部分开销。第6章同时展示准备、后续解密和保留叶钥大小，以判断实际使用方式。

\section{本章小结}

本章说明了公共接口背后的具体计算。LATTE通过NTRU结构陷门完成两层委派，Lai通过追加身份块
形成叶陷门；BNN转换、Luo、SRGG、Sageloli和Wang--PW分别提供不同结构的身份签名。六种IBS
在公共协议中的职责相同，但签名生成、对象尺寸和运行资源差异明显。

三项优化直接对应这些差异。Sageloli利用环结构把大分解转成逐频率低秩计算，SRGG按字段范围
编码四类对象，Lai会话复用固定身份对的叶密钥。第~6章将用实际时间、尺寸和内存说明每项方法
带来的变化。
